%% ============================================================================
%% revised_abstract.tex
%%
%% Replacement for the abstract environment of paper.tex (original lines
%% 101--142). Calibrated novelty per referee report:
%%   - foregrounds the genuinely new content
%%       (i)   chi(c) = chi_w(c^q) factorisation,
%%       (ii)  the 23-/39-class stratified atlas,
%%       (iii) the mu_q / mu_{q'} symmetry bookkeeping,
%%       (iv)  the p-visibility statement (AT7) as the one new analytic
%%             observation;
%%   - identifies chi_w as the classical edge-indicial polynomial in w = c^q;
%%   - downplays AT1/AT3/AT5/AT7 as analytic theorems and labels their
%%     classical content explicitly (Levelt--Turrittin, Sibuya, Wasow);
%%   - removes overclaims about wild character varieties;
%%   - states clearly what is proved vs. classical.
%%
%% Word count: 226.
%% ============================================================================
\begin{abstract}
For a scalar local linear operator $L$ (differential, difference, or
$q$-difference) with one irregular singularity at $x = \infty$ whose
Newton polygon has a unique dominant edge of slope $-p/q$ with $r{+}1$
lattice points, the WKB ansatz on that edge yields a characteristic
polynomial $\chi \in \mathbb{C}[c]$. Our central algebraic identity is
\[
   \chi(c) \;=\; \chi_w(c^q),
\]
where $\chi_w \in \mathbb{C}[w]$ is the classical edge-indicial
polynomial of degree $r$ in the variable $w = c^q$. The exponent
constellation $\Cl(L) \subset \mathbb{C}$ therefore decomposes as a
$\mu_q$-equivariant union of $r$ fibres of $c \mapsto c^q$ over
$\Roots(\chi_w)$.

This factorisation organises a finite stratified atlas. A computational
sweep produces $23$ constellation classes for single-edge slopes
$-p/q \in \{1/2,\,1/3,\,1/4,\,2/3,\,3/2\}$ and rank $r \le 3$,
refining to $39$ classes once multi-edge polygons, nested ramification,
and the prime-$q$ arithmetic stratification are tracked. We prove a
per-edge factorisation along the slope filtration, an iterated
$\mu_{q'}$-pullback law for nested inner polynomials, a block-diagonal
pullback for systems, and a prime-$q$ Kummer dichotomy separating the
algebraically independent stratum from the cyclotomic stratum.

The analytic input attaching anti-Stokes directions to the constellation
is classical (Levelt--Turrittin, Sibuya, Wasow); on the $q$-fold cover
the rays for a generic single-edge operator satisfy $\arg \Dij + p\phi
\equiv \pi/2 \pmod{\pi}$. From this we extract one explicit consequence:
the slope numerator $p$, invisible to $\Cl(L)$ itself, is visible to
the cover-side ray pattern.

The novel content of the paper is the algebraic identity, the
stratified atlas, the $\mu_q/\mu_{q'}$ symmetry bookkeeping, and the
$p$-visibility observation. The analytic statements (\Cref{thm:AT1},
\Cref{thm:AT3}, \Cref{thm:AT5}, \Cref{thm:AT7}) are explicit
$\mu_q/\mu_{q'}$-equivariant repackagings of classical formal-classification
theory; we make no claim of new wild character variety theorems.
\end{abstract}
